% SOFTWARE_TEST_PLAN.tex
% Software Test Plan for the AWEB toolchain (atangle/aweave) release 2.2.
% Conforms to IEEE Std 829-2008 (Standard for Software and System Test
% Documentation).  This document is intended to be typeset with pdflatex
% but uses only the standard article class so that any TeX engine of
% LaTeX2e vintage will produce the document.

\documentclass[11pt]{article}

\usepackage[margin=1in]{geometry}
\usepackage{tabularx}
\usepackage{longtable}
\usepackage{booktabs}
\usepackage{array}
\usepackage{enumitem}
\usepackage{hyperref}
\usepackage{parskip}

\hypersetup{
  pdftitle  = {AWEB-STP-2.2: Software Test Plan},
  pdfauthor = {AWEB Toolchain Maintainers},
  colorlinks = true,
  linkcolor  = blue,
  urlcolor   = blue
}

\newcommand{\stpid}{AWEB-STP-2.2}
\newcommand{\sut}{\textsf{atangle}/\textsf{aweave}}

\title{Software Test Plan\\\large AWEB Toolchain (\textsf{atangle} 2.2 / \textsf{aweave} 2.2)}
\author{AWEB Toolchain Maintainers}
\date{2026--05--06}

\begin{document}
\maketitle

\begin{center}
\begin{tabular}{ll}
\toprule
Document identifier   & \stpid\\
Status                & Approved for Test\\
Conforming standard   & IEEE Std 829--2008\\
\bottomrule
\end{tabular}
\end{center}

\tableofcontents
\bigskip

\section{Introduction}

\subsection{Document identifier}
This is the Software Test Plan (STP) for release 2.2 of the AWEB
toolchain, comprising the programs \texttt{atangle} and \texttt{aweave}.
It is identified as \stpid.

\subsection{Scope}
Release~2.0 extended the toolchain so that AWEB source files may contain
Ada~95 (ISO/IEC~8652:1995) programs in addition to the previously
supported Ada~83 (ANSI/MIL-STD-1815A-1983) programs.  Six reserved
words were added to the lexical recognizer of \texttt{aweave}:
\begin{quote}
\texttt{abstract},\ \texttt{aliased},\ \texttt{protected},\
\texttt{requeue},\ \texttt{tagged},\ \texttt{until}.
\end{quote}
Release~2.1 added a \texttt{-v} / \texttt{--version} command-line option
to both programs and removed the bundled PostgreSQL stack.

Release~2.2 extends the toolchain to recognise programs written in
Ada~2012 (ISO/IEC~8652:2012) in addition to Ada~83 and Ada~95.  Four
reserved words --- three carried over from Ada~2005 and one introduced
by Ada~2012 itself --- are added to the lexical recognizer of
\texttt{aweave}:
\begin{quote}
\texttt{interface},\ \texttt{overriding},\ \texttt{synchronized}\
(Ada~2005);\ \texttt{some}\ (Ada~2012).
\end{quote}
Because the lexical structure of Ada~2012 is unchanged with respect to
identifiers, \texttt{atangle} requires no syntactic change; the new
keywords are passed through verbatim and pretty-printed by
\texttt{aweave}.  The container is now tagged
\texttt{evansjr/aweb:2.2}, and the version banner of both programs has
been updated to ``Version 2.2 (Ada 2012)''.

The change-file workflow (requirement~11) and the \texttt{-v} /
\texttt{--version} command-line flag (requirement~8) are carried over
unchanged from release~2.1; only the banner strings they print have
been retargeted to the new version.

This STP exercises the unchanged behaviour (regression), the Ada~95
keyword classification (regression), the new Ada~2012 keyword
classification, the version-flag behaviour, and the surrounding
container artefacts (\texttt{Dockerfile}, \texttt{aweb.sh},
\texttt{docker-entrypoint.sh}).

\subsection{References}
\begin{enumerate}[label=R\arabic*., leftmargin=2.5em]
  \item IEEE Std~829--2008, \emph{Standard for Software and System Test
        Documentation}.
  \item ISO/IEC~8652:2012, \emph{Information technology --- Programming
        languages --- Ada} (Ada~2012).
  \item ISO/IEC~8652:1995, \emph{Information technology --- Programming
        languages --- Ada} (Ada~95).
  \item ANSI/MIL-STD-1815A-1983, \emph{Reference Manual for the Ada
        Programming Language} (Ada~83).
  \item D.~E.~Knuth, \emph{Literate Programming}, CSLI, 1992.
  \item \texttt{requirements.txt} in this repository: functional
        requirements for release~2.1.
  \item \texttt{requirements2.txt} in this repository: functional
        requirements for release~2.2.
  \item \texttt{aweave/aweave.aweb}, \texttt{atangle/atangle.aweb}:
        literate sources of the SUT.
\end{enumerate}

\section{Test items}

\begin{center}
\begin{tabular}{@{}llll@{}}
\toprule
ID    & Item                          & Version & Source of record\\
\midrule
TI--1 & \texttt{atangle} executable   & 2.2 & \texttt{atangle/atangle.adb} (+ tangled units)\\
TI--2 & \texttt{aweave}  executable   & 2.2 & \texttt{aweave/aweave.adb}  (+ tangled units)\\
TI--3 & \texttt{Dockerfile}           & 2.2 & \texttt{./Dockerfile}\\
TI--4 & \texttt{aweb.sh} driver     & 2.2 & \texttt{./aweb.sh}\\
TI--5 & \texttt{docker-entrypoint.sh} & 2.2 & \texttt{./docker-entrypoint.sh}\\
\bottomrule
\end{tabular}
\end{center}
Each item is delivered as part of a single OCI image tagged
\texttt{evansjr/aweb:2.2}.

\section{Features to be tested}

\begin{center}
\begin{tabularx}{\linewidth}{@{}lX@{}}
\toprule
ID   & Feature\\
\midrule
F--1  & Build of \texttt{atangle} and \texttt{aweave} from delivered Ada
        source (\texttt{gnatmake}).\\
F--2  & Banner / version reporting.\\
F--3  & Tangling of an Ada~83 program (regression).\\
F--4  & Tangling of an Ada~95 program containing the six new reserved
        words.\\
F--5  & Weaving of an Ada~83 program (regression).\\
F--6  & Weaving of an Ada~95 program: each new reserved word is set in
        \verb|\&{...}| (bold) face.\\
F--6b & Weaving of an Ada~2012 program: the four new Ada~2005 / Ada~2012
        reserved words (\texttt{interface}, \texttt{overriding},
        \texttt{synchronized}, \texttt{some}) are set in
        \verb|\&{...}| (bold) face.\\
F--7  & Round-trip: tangling is idempotent on a fixed Ada~95 source.\\
F--8  & Container build (\texttt{docker build}) yields image
        \texttt{evansjr/aweb:2.2}.\\
F--9  & \emph{Withdrawn in release 2.1}: PostgreSQL was removed from the
        image (requirement~10), so the previous initialisation /
        persistence feature no longer applies.\\
F--10 & \texttt{make}, \texttt{file}, \texttt{emacs} are present in the
        container image.\\
F--11 & \texttt{aweb.sh} dispatch (\texttt{atangle}, \texttt{aweave},
        \texttt{shell}, \texttt{test}) functions inside the container.\\
F--12 & The full test harness can be executed inside the container
        (requirement~7).\\
F--13 & The \texttt{-v} / \texttt{--version} command-line flag prints the
        version banner and exits successfully without opening any input
        or output files (requirement~8).\\
F--14 & Running \texttt{atangle} on \texttt{atangle.aweb}/\texttt{aweave.aweb}
        together with the bundled \texttt{.ach} change file produces the
        GNAT-style \texttt{.adb}/\texttt{.ads} output files
        (requirement~11).\\
\bottomrule
\end{tabularx}
\end{center}

\section{Features not to be tested}

\begin{center}
\begin{tabularx}{\linewidth}{@{}lXX@{}}
\toprule
ID    & Feature & Rationale\\
\midrule
NF--1 & Ada~2022 keywords (\texttt{parallel}; \texttt{declare} as
        expression form)
      & Out of scope for release~2.2; only Ada~2012 was contracted.\\
NF--2 & Performance / throughput on multi-megabyte AWEB inputs
      & No performance requirement was levied.\\
NF--3 & Internationalization of TeX output
      & English/ASCII output only is contracted.\\
NF--4 & TeX rendering pipeline (\texttt{tex}, \texttt{pdflatex})
      & Downstream of the SUT; the STP verifies textual output, not
        typesetting.\\
\bottomrule
\end{tabularx}
\end{center}

\section{Approach}
The strategy is \emph{black-box functional testing} driven by a shell
harness (\texttt{tests/test\_aweb.sh}) that
\begin{enumerate}[leftmargin=*]
  \item compiles both programs with \texttt{gnatmake} (or detects an
        existing in-container installation under
        \texttt{/usr/local/bin}),
  \item runs each program against curated AWEB fixtures and inspects
        the output for keyword classification and structural
        properties,
  \item performs a round-trip equivalence test for the new keywords,
  \item inspects the container image for tooling and version labels.
\end{enumerate}
Each test case is binary pass/fail.  The harness exits with status~0
only if every case passes; it prints a one-line summary per case.  It is
suitable for CI pipelines and for ad-hoc developer verification.  No
interactive judgement is required.

\subsection{Test levels covered}
\begin{description}[leftmargin=*]
  \item[Unit] (lexical-table) F--2, F--4, F--6, F--6b.
  \item[Integration] (program + I/O) F--1, F--3, F--5, F--7, F--11.
  \item[System] (container) F--8, F--9, F--10, F--12.
\end{description}

\subsection{Test types}
Functional, regression, configuration.

\subsection{Execution environments}
The harness is designed to run in two execution environments:
\begin{enumerate}[leftmargin=*]
  \item \emph{Host} --- requires Docker for the container-level cases;
        requires GNAT for the native-compile cases (when not running in
        the container).
  \item \emph{In container} --- invoked via
        \texttt{aweb.sh test} or \texttt{aweb.sh shell} followed by
        \texttt{bash tests/test\_aweb.sh}.  In this mode the harness
        uses the pre-installed binaries at \texttt{/usr/local/bin/atangle}
        and \texttt{/usr/local/bin/aweave} and skips the docker-of-docker
        cases (TC--010 \dots TC--014).
\end{enumerate}

\section{Item pass/fail criteria}
A test item passes when \emph{every} test case mapped to it (see~\S\ref{sec:cases})
passes.  A test case passes when its observed result matches its
expected result exactly: \texttt{diff~-u} empty for textual cases, exit
status zero for build cases, presence of artefact for configuration
cases.

\section{Suspension criteria and resumption requirements}
Testing is suspended if
\begin{itemize}[leftmargin=*]
  \item \texttt{gnatmake} fails on unmodified fixtures (toolchain rather
        than SUT failure),
  \item the Docker daemon is unreachable when running container-level
        cases on the host, or
  \item a fixture file in \texttt{tests/} is reported missing.
\end{itemize}
Testing resumes once the blocking environmental fault is corrected; all
previously executed cases shall be re-executed.

\section{Test deliverables}

\begin{center}
\begin{tabularx}{\linewidth}{@{}lX@{}}
\toprule
Deliverable                        & Location\\
\midrule
This Software Test Plan            & \texttt{SOFTWARE\_TEST\_PLAN.tex}\\
Test harness                       & \texttt{tests/test\_aweb.sh}\\
Test fixtures (Ada~83, Ada~95, Ada~2012) & \texttt{tests/fixtures/}\\
Test log (produced by harness)     & \texttt{tests/test\_aweb.log}\\
\bottomrule
\end{tabularx}
\end{center}

\section{Testing tasks}

\begin{center}
\begin{tabular}{@{}rll@{}}
\toprule
\# & Task & Predecessors\\
\midrule
1 & Stage source tree                                   & ---\\
2 & Build container image                               & 1\\
3 & Run \texttt{aweb.sh test} (or run the harness on the host) & 1, 2\\
4 & Inspect log; declare pass/fail                      & 3\\
\bottomrule
\end{tabular}
\end{center}

\section{Test cases}\label{sec:cases}

The fifteen cases below are implemented by
\texttt{tests/test\_aweb.sh}.  ``Procedure'' is summarised; the full
procedure is the script itself.

{\small
\begin{longtable}{@{}lp{3.5cm}lp{4.6cm}p{4.0cm}@{}}
\toprule
ID & Title & Maps to & Procedure (abbreviated) & Expected result\\
\midrule
\endhead
TC--001 & Build atangle & F--1 &
\texttt{cd atangle \&\& gnatmake atangle.adb} &
Exit~0; \texttt{atangle} executable produced.\\

TC--002 & Build aweave & F--1 &
\texttt{cd aweave \&\& gnatmake aweave.adb} &
Exit~0; \texttt{aweave} executable produced.\\

TC--003 & atangle banner & F--2 &
\texttt{atangle ada83.aweb} ; first stdout line. &
Contains ``This is ATANGLE, Version 2.2 (Ada 2012)''.\\

TC--004 & aweave banner & F--2 &
\texttt{aweave ada83.aweb} ; first stdout line. &
Contains ``This is AWEAVE, Version 2.2 (Ada 2012 keyword support)''.\\

TC--005 & Tangle Ada 83 fixture & F--3 &
Tangle \texttt{ada83.aweb}; check produced \texttt{.a} file. &
Output produced; no errors.\\

TC--006 & Tangle Ada 95 fixture & F--4 &
Tangle \texttt{ada95.aweb}; check produced \texttt{.a} file. &
Output produced; no errors.\\

TC--007 & Weave Ada 83 fixture (regression) & F--5 &
Weave \texttt{ada83.aweb}; check that
\verb|\&{begin}|, \verb|\&{end}| appear in the \texttt{.tex}. &
Both control sequences present.\\

TC--008 & Weave Ada 95: new keywords are bold & F--6 &
Weave \texttt{ada95.aweb}; verify
\verb|\&{abstract}|, \verb|\&{aliased}|, \verb|\&{protected}|,
\verb|\&{requeue}|, \verb|\&{tagged}|, \verb|\&{until}| each
appear at least once in the \texttt{.tex}. &
All six present.\\

TC--008b & Weave Ada 2012: new keywords are bold & F--6b &
Weave \texttt{ada2012.aweb}; verify
\verb|\&{interface}|, \verb|\&{overriding}|,
\verb|\&{synchronized}|, \verb|\&{some}| each
appear at least once in the \texttt{.tex}. &
All four present.\\

TC--009 & Round-trip equivalence & F--7 &
Tangle \texttt{ada95.aweb} twice; \texttt{diff} the two outputs. &
\texttt{diff} empty.\\

TC--010 & Container build & F--8 &
\texttt{docker build -t evansjr/aweb:2.2 .} &
Build succeeds; image carries label
\texttt{org.opencontainers.image.version=2.2}.\\

TC--011 & Container has make, file, emacs & F--10 &
\texttt{docker run \dots\ command -v make file emacs}. &
All three resolve to absolute paths.\\

TC--012 & \multicolumn{4}{l}{\emph{Withdrawn in release 2.1: PostgreSQL
was removed from the image per requirement~10.}}\\

TC--013 & \texttt{aweb.sh} dispatch & F--11 &
\texttt{aweb.sh atangle ada95.aweb} inside the container. &
Exit~0; \texttt{web\_output.a} produced.\\

TC--014 & In-container Ada~95 keyword classification & F--6, F--12 &
Run \texttt{aweave ada95.aweb} in the container; check the produced
\texttt{.tex} contains \verb|\&{abstract}|, \verb|\&{aliased}|,
\verb|\&{protected}|, \verb|\&{requeue}|, \verb|\&{tagged}|,
\verb|\&{until}|. &
All six classified as reserved (independent of host GNAT).\\

TC--015 & \texttt{atangle -v} prints version & F--13 &
Run \texttt{atangle -v} in an empty working directory; observe
exit~0, stdout contains the banner, and no \texttt{.a}/\texttt{.tex}
files are produced. &
Banner contains
``This is ATANGLE, Version 2.2 (Ada 2012)''; no artefacts.\\

TC--016 & \texttt{aweave --version} prints version & F--13 &
Run \texttt{aweave --version} in an empty working directory; observe
exit~0, stdout contains the banner, and no \texttt{.a}/\texttt{.tex}
files are produced. &
Banner contains
``This is AWEAVE, Version 2.2 (Ada 2012 keyword support)'';
no artefacts.\\

TC--017 & Retangle \texttt{atangle.aweb} with \texttt{atangle.ach} & F--14 &
Copy \texttt{atangle.aweb} and \texttt{atangle.ach} into a clean
directory, run \texttt{atangle atangle.aweb atangle.ach}, and inspect
the produced files. &
\texttt{atangle.adb}, \texttt{data\_structures.\{adb,ads\}},
\texttt{hashing.\{adb,ads\}}, \texttt{input\_output.\{adb,ads\}},
\texttt{input\_phase.\{adb,ads\}}, \texttt{int\_number\_io.ads},
\texttt{output\_phase.\{adb,ads\}} all present; no errors reported.\\

TC--018 & Retangle \texttt{aweave.aweb} with \texttt{aweave.ach} & F--14 &
Copy \texttt{aweave.aweb} and \texttt{aweave.ach} into a clean
directory, run \texttt{atangle aweave.aweb aweave.ach}, and inspect
the produced files. &
\texttt{aweave.adb}, \texttt{data\_structures.\{adb,ads\}},
\texttt{first\_phase.\{adb,ads\}},
\texttt{input\_output.\{adb,ads\}},
\texttt{lexical\_scanning.\{adb,ads\}},
\texttt{output.\{adb,ads\}},
\texttt{parsing.\{adb,ads\}},
\texttt{parsing-ada\_parse.adb}, \texttt{parsing-do\_cases.adb},
\texttt{parsing-translate.adb},
\texttt{searching.\{adb,ads\}},
\texttt{second\_phase.\{adb,ads\}},
\texttt{third\_phase.\{adb,ads\}} all present; no errors reported.\\
\bottomrule
\end{longtable}
}

\section{Environmental needs}
\begin{itemize}[leftmargin=*]
  \item Host with Docker $\geq$~24 and $\geq$~2~GiB RAM.
  \item GNAT (any FSF release with GCC~$\geq$~9) for non-container runs
        of TC--001\dots TC--009.  Inside the container GNAT is provided
        by the image and the user need not install anything on the
        host.
  \item Coreutils (\texttt{diff}, \texttt{grep}, \texttt{sed},
        \texttt{cmp}) for the harness.
\end{itemize}

\section{Responsibilities}

\begin{center}
\begin{tabular}{@{}ll@{}}
\toprule
Role               & Party\\
\midrule
Test author        & Toolchain maintainer\\
Test executor      & CI system or developer running
                     \texttt{aweb.sh test}\\
Defect adjudicator & Toolchain maintainer\\
\bottomrule
\end{tabular}
\end{center}

\section{Staffing and training}
No specialised training is required beyond the ability to read AWEB
source, Ada source, and shell scripts.

\section{Schedule}
The harness is self-contained and intended to run on every commit; a
full pass typically completes in under thirty seconds on a developer
workstation, plus the time to build the container image
($\approx$~2~minutes on first build, seconds when cached).

\section{Risks and contingencies}

\begin{center}
\begin{tabularx}{\linewidth}{@{}XX@{}}
\toprule
Risk & Mitigation\\
\midrule
Future GNAT release reserves additional Ada~2022 keywords used as
identifiers in legacy AWEB source. &
Restrict regression fixtures to identifiers known not to clash; pin the
GNAT version in the \texttt{Dockerfile}.\\
Diff-based weaving comparison too brittle (whitespace).
& Compare on a control-sequence whitelist instead of full text.\\
Docker daemon unavailable in CI environment.
& Container-level cases (TC--010\dots TC--014) skip with explicit
\texttt{SKIP} status; the harness still asserts non-container cases
pass.\\
\bottomrule
\end{tabularx}
\end{center}

\section{Approvals}

\begin{center}
\begin{tabular}{@{}lll@{}}
\toprule
Role               & Name             & Date\\
\midrule
Author             & (toolchain team) & 2026-05-06\\
Configuration mgr. & (toolchain team) & 2026-05-06\\
Quality assurance  & (toolchain team) & 2026-05-06\\
\bottomrule
\end{tabular}
\end{center}

\end{document}
